\section{本章小结}

本章针对有限元空间离散过程中由畸变单元引起的局部高频人造刚度与全局低频响应之间的软硬耦合问题，提出并测试了一种基于算子感知的局部基空间重构与隔离方法。为了缓解畸变单元对全局系统矩阵条件数的影响，本章在不改变原始网格宏观节点连接拓扑的前提下，通过对退化单元区域实施局部多面体聚合，并在这些聚合区域内引入狄拉克小波（Dirac-wavelet）转换策略，重构了局部的离散解空间。该机制对由畸变网格引入的高频奇异刚度成分进行了隔离处理，同时保留了系统矩阵的稀疏构造。数值实验表明，本章方法降低了含低质量单元系统的最大特征值，将预处理共轭梯度算法的求解效率提升了约 $20\%$，并在涉及切割破坏等动态仿真场景中保持了数值稳定性。就本章范围而言，这说明对于由畸变单元引起的离散化软硬耦合问题，可以通过局部表示空间的重构与隔离来削弱其对主求解过程的数值干扰。

在全文脉络中，本章对应的是一条区别于第二章的处理路径。第二章通过全局广义坐标变换来消除链式结构中的拉伸惩罚能量；本章所面对的则是三维非结构化网格中的高频连接模式，这类模式更局部，也更难通过统一的全局运动学重构来处理。因此，本章采用局域化的表示空间微调策略：通过识别并隔离那些在宏观物理尺度上影响较小、但会干扰底层代数求解的高频奇异基底，减弱了局部人造刚度对主求解空间的数值干扰。

然而，第二章的全局运动学链式重构与本章的局部基空间隔离，处理的仍然都是局部高频刚度对主求解空间的数值干扰。在另一类典型场景中，病态性的主要来源并不直接来自显式几何约束或局部畸变单元，而是来自降维特征空间本身的构造方式。在分布式区域分解框架下，传统界面缩减方法受限于固定界面刚度，较难充分捕获全局低频“软”响应。如何重构降维特征子空间，缓解固定划分与截断误差带来的表达不足，是下一章将要讨论的问题。
